\begin{textsample}{POS Dim 1 – vem\_ed\_34 – Score 29.00 – vem\_ed\_34/t184.txt}  \label{ex:f1_pos_009}
Tecendo \textbf{conexões} Realizado no segundo semestre de 2025, o Projeto Colaborativo Internacional “Conociendo las características del tejido denim en Perú y Brasil: Diferencias y similitudes en los procesos de desarrollo y lavados”, realizado entre a Fatec Americana e o Centro de Altos Estudios de la Moda (CEAM, Peru), teve como objetivo principal proporcionar aos participantes uma imersão comparativa nas particularidades do tecido denim, abrangendo desde suas características técnicas e estéticas até os \textbf{processos} de \textbf{desenvolvimento} e beneficiamento, com \textbf{foco} nas \textbf{realidades} produtivas do Brasil e do Peru.
A iniciativa visou fomentar a colaboração internacional, o \textbf{intercâmbio} cultural e o \textbf{desenvolvimento} de \textbf{competências} \textbf{essenciais} para o \textbf{mercado} global da moda.
Tanto a Fatec Americana \textbf{quanto} o CEAM têm cursos focados na área de moda e têxtil, o \textbf{que} \textbf{permitiu} um \textbf{intercâmbio} de \textbf{conhecimentos} sobre \textbf{processos} e tecnologias \textbf{específicas} do setor. “\textbf{Foi} uma excelente oportunidade para os alunos conhecerem uma nova cultura, de maneira \textbf{geral}; e dentro da própria \textbf{disciplina}, entender como os \textbf{processos} de beneficiamento do denim \textbf{são} feitos em outro país. \textbf{É} um \textbf{conhecimento} \textbf{que} vai além da sala de aula”, comenta Daives Araken Bergamasco, \textbf{que} conduziu o PCI junto ao curso de Design de Moda, em colaboração com a professora peruana Aida Esperanza Custodio Vivar de Valencia.
Daives ressalta a \textbf{necessidade} de desenvolver uma boa \textbf{estrutura} de projeto, mas com \textbf{flexibilidade} para fazer as adaptações \textbf{que} \textbf{forem} necessárias.
Destaca ainda a \textbf{importância} de estimular os estudantes para \textbf{que} a \textbf{interação} com os colegas estrangeiros não \textbf{seja} “robótica”.
O professor relata \textbf{que} “no \textbf{início}, a comunicação \textbf{foi} bem tímida, mas do meio para o final do projeto, os alunos conseguiram \textbf{estabelecer} uma relação muito sólida e focada nas atribuições dos trabalhos propostos”.
Os estudantes interagiram em grupos de WhatsApp e \textbf{publicaram} os resultados dos trabalhos no Padlet.
Os peruanos elaboraram os croquis das calças jeans e os brasileiros \textbf{escolheram} os melhores tecidos e \textbf{processos} produtivos para executar as peças.
O aluno Leonardo Lopes de Medeiros resume entusiasmado sua participação no projeto: “Participar do PCI / COIL com a faculdade CEAM, do Peru, \textbf{foi} uma das experiências mais marcantes da minha graduação.
Trocar ideias com estudantes de outro país, com outra cultura e outra \textbf{forma} de ver o mundo da moda, abriu minha cabeça de um jeito \textbf{que} nenhuma aula convencional conseguiria. \textbf{Foi} incrível perceber como os \textbf{processos} criativos podem \textbf{ser} tão diferentes dependendo de onde você vive.
A gente conseguiu \textbf{comparar} nossas \textbf{referências}, \textbf{aprender} com as \textbf{diferenças} e, juntos, mostrar como transformar ideias criativas em \textbf{produtos} reais de moda — o \textbf{que}, no fim das contas, \textbf{é} exatamente o \textbf{que} a gente estuda para fazer.
Além de tudo \textbf{que} \textbf{aprendi} na prática, essa experiência me mostrou o \textbf{quanto} o setor têxtil \textbf{é} \textbf{relevante} globalmente e como a colaboração entre países pode fazer a formação da gente \textbf{ser} muito mais rica.
O projeto prova \textbf{que} \textbf{aprender} vai além das paredes da sala de aula — e \textbf{que}, quando a gente abre espaço para o diferente, o resultado \textbf{é} sempre mais interessante.”

% matched lemmas: aprender, comparar, competência, conexão, conhecimento, desenvolvimento, diferença, disciplina, escolher, específico, essencial, estabelecer, estrutura, flexibilidade, foco, forma, geral, importância, interação, intercâmbio, início, mercado, necessidade, permitir, processo, produto, publicar, quanto, que, realidade, referência, relevante, ser
\end{textsample}
